% =====================================================================
%  THÈME 2 — Notation scientifique & unités — CORRIGÉ — Niveau FACILE
% =====================================================================
\documentclass[11pt,a4paper]{article}
\usepackage[utf8]{inputenc}
\usepackage[T1]{fontenc}
\usepackage{lmodern}
\usepackage[french]{babel}
\usepackage{amsmath,amssymb,mathtools}
\usepackage{icomma}
\usepackage{siunitx}
\sisetup{output-decimal-marker={,},per-mode=symbol}
\usepackage[version=4]{mhchem}
\usepackage{xcolor}
\usepackage[most]{tcolorbox}
\usepackage{enumitem}
\usepackage{fancyhdr}
\usepackage[a4paper,margin=2.2cm,headheight=26pt,headsep=8pt]{geometry}

\definecolor{primary}{RGB}{150,98,162}
\definecolor{primarydark}{RGB}{92,52,110}
\newcommand{\cs}{chiffre significatif}
\newcommand{\CS}{chiffres significatifs}

\pagestyle{fancy}\fancyhf{}
\fancyhead[L]{\small\textcolor{primarydark}{\textbf{Fiche 1} — Chiffres significatifs}}
\fancyhead[R]{\small\textcolor{primarydark}{\textsc{Corrigé · Facile · Thème 2}}}
\fancyfoot[C]{\small\thepage}
\renewcommand{\headrulewidth}{0.8pt}
\renewcommand{\headrule}{\hbox to\headwidth{\color{primary}\leaders\hrule height \headrulewidth\hfill}}

\newtcolorbox[auto counter]{cor}[1][]{enhanced,breakable,boxrule=0.8pt,arc=2pt,
  colback=primarydark!5,colframe=primarydark,left=3mm,right=3mm,top=2mm,bottom=2mm,
  fonttitle=\bfseries,coltitle=white,
  title={Correction — Exercice~\thetcbcounter\ifx\relax#1\relax\relax\else\ ---\ \textmd{\itshape #1}\fi}}

\setlength{\parindent}{0pt}\setlength{\parskip}{3pt}

\begin{document}

\begin{center}
{\Large\bfseries\color{primarydark} Corrigé — Niveau Facile}\\[1pt]
{\color{primary} Thème 2 : notation scientifique \& changements d'unité}
\end{center}
\vspace{2mm}

\begin{cor}[Grands nombres $\rightarrow$ notation scientifique]
On recule la virgule jusqu'après le premier chiffre non nul ; $n$ est positif.
\begin{itemize}[leftmargin=1.6em,itemsep=1pt]
  \item a) $\num{45000}=\num{4,5e4}$ \quad b) $\num{128}=\num{1,28e2}$ \quad c) $\num{9300}=\num{9,3e3}$
  \item d) $\num{6}=\num{6e0}$ \quad e) $\num{20960}=\num{2,0960e4}$ (les $4$~CS sont conservés).
\end{itemize}
\end{cor}

\begin{cor}[Petits nombres $\rightarrow$ notation scientifique]
On avance la virgule jusqu'après le premier chiffre non nul ; $n$ est négatif.
\begin{itemize}[leftmargin=1.6em,itemsep=1pt]
  \item a) $\num{0,02096}=\num{2,096e-2}$ \quad b) $\num{0,00072}=\num{7,2e-4}$ \quad c) $\num{0,5}=\num{5e-1}$
  \item d) $\num{0,0000130}=\num{1,30e-5}$ (le zéro de queue reste !) \quad e) $\num{0,048}=\num{4,8e-2}$
\end{itemize}
\end{cor}

\begin{cor}[Notation scientifique $\rightarrow$ écriture décimale]
\begin{itemize}[leftmargin=1.6em,itemsep=1pt]
  \item a) $\num{2,096e-2}=\num{0,02096}$ \quad b) $\num{3,00e4}=\num{30000}$ \quad c) $\num{6,02e2}=\num{602}$
  \item d) $\num{1,0e-7}=\num{0,00000010}$ \quad e) $\num{4,5e1}=\num{45}$
\end{itemize}
\end{cor}

\begin{cor}[Ne garder que les chiffres significatifs dans la mantisse]
\begin{itemize}[leftmargin=1.6em,itemsep=1pt]
  \item a) $\num{0,004500}$ : $4$~CS $\rightarrow \num{4,500e-3}$ (on garde les deux zéros de queue).
  \item b) $\num{0,0100}$ : $3$~CS $\rightarrow \num{1,00e-2}$.
  \item c) $\num{15,200}$ : $5$~CS $\rightarrow \num{1,5200e1}$.
\end{itemize}
Dans les trois cas, la mantisse contient exactement les CS du nombre de départ.
\end{cor}

\begin{cor}[Vrai ou faux : est-ce une notation scientifique correcte ?]
La condition est $1\le|a|<10$.
\begin{itemize}[leftmargin=1.6em,itemsep=1pt]
  \item a) $\num{12,5e3}$ : \textbf{faux} ($a=\num{12,5}\ge10$). Correct : $\num{1,25e4}$.
  \item b) $\num{3,2e5}$ : \textbf{correct}.
  \item c) $\num{0,4e-2}$ : \textbf{faux} ($a=\num{0,4}<1$). Correct : $\num{4e-3}$.
  \item d) $\num{8,00e0}$ : \textbf{correct} (vaut $\num{8,00}$, $3$~CS ; $10^0=1$ est admis).
\end{itemize}
\end{cor}

\end{document}
